\documentclass[10pt,a4paper]{amsart}
\usepackage[margin=19mm]{geometry}
\usepackage{amsmath,amssymb,mathtools}
\usepackage[scr=rsfs,cal=euler]{mathalfa}
\usepackage{xfrac}
\usepackage[unicode,hidelinks,bookmarks=false]{hyperref}
\newcommand{\ie}{i.e.\ }
\newcommand{\cf}{cf.\ }
\newcommand{\resp}{respectively\ }
\newcommand{\Subsection}{\subsection}
\providecommand{\nobreakdash}{\nobreak\hbox{-}}
\setlength{\emergencystretch}{2em}
\multlinegap0pt
\usepackage{amssymb}
\usepackage[shortlabels]{enumitem}
\newtheorem{lemma}{Lemma}
\newcommand{\F}{\mathbf{F}}
\renewcommand{\Re}{\op{Re}}
\newcommand\br[1]{{\left(#1\right)}}
\newcommand{\gen}[1]{\langle{#1}\rangle}
\def\one{\mathbf{1}}
\newcommand{\op}{\operatorname}
\title{Expanding groups with large diameter}
\author{Sean Eberhard, Luca Sabatini}
\date{Source: arXiv:2602.13582v1 (CC BY 4.0)}
\begin{document}
\maketitle

\noindent\emph{Source and licence.} Expanding groups with large diameter. Version arXiv:2602.13582v1 (CC BY 4.0), distributed by arXiv under the Creative Commons Attribution 4.0 International licence, http://creativecommons.org/licenses/by/4.0/, as recorded in the arXiv licence metadata for this exact version and shown on the abstract page https://arxiv.org/abs/2602.13582v1. Full licence notice: \texttt{licenses/arxiv-2602.13582-CC-BY-4.0.txt}.\par\smallskip

\noindent\textit{Editorial scope.} Counted unit: the article's lemma lem:switching (source0.tex lines 468--476). The article's complete original proof of it is delivered with it (source0.tex lines 477--521, one \texttt{proof} environment of 45 lines). No mathematics is written, repaired or re-derived: the statement and the proof are the article's own lines, in the article's own notation, and no retained line is substituted or re-flowed.  No further source-local context is needed: every cross-reference of every retained line resolves inside the two delivered spans.  Nothing was omitted from any retained span, so the package carries no omission record.  No retained line contains a citation, so the wrapper carries no bibliography and no bibliography entry.  No retained line contains a figure environment, an image-inclusion command or drawing code, so no image is delivered and the package declares no asset dependency.  Editorial formatting only, in addition: an AMS \texttt{amsart} preamble that restates the article's own declarations and macros used by the retained lines, namely the environments \texttt{lemma} on separate counters, together with the standard packages those lines need.  The retained environments print in this document as Lemma 1 in order of appearance, whereas the article prints the same items with the numbering of its own full document.  The complete native source of the article as distributed in the arXiv e-print for this exact version is bundled unchanged as \texttt{sources/194654/source0.tex}.
\begin{lemma}
    \label{lem:switching}
    Let $v \in V_0$. We have
    \[
       \max_{w \in V \setminus \gen\one} |\lambda_{v, w}|^2
       \> \le \>
       \frac12 + \frac12 \max_{0 \ne u \in \F_p} |\lambda_v(u)|^2.
    \]
\end{lemma}

\begin{proof}
    We apply a switching argument based on the Cauchy--Schwarz inequality.

    Let $w \in V \setminus \gen\one$ and let us bound $\lambda_{v,w}$.
    Obviously $\lambda_{v, w} = \lambda_{v, w^\sigma}$ for all $\sigma \in S_n$, so we may assume $w_1 \ne w_2$ without loss of generality. Let $u = w_1 - w_2$ and let $\tau$ be the transposition $(1,2)$. Then
    \[
        \lambda_{v,w}
        = \frac1{n!} \sum_{\sigma \in S_n} e_p( \gen{v, w^\sigma})
        = \frac1{n!} \sum_{\sigma \in S_n} \frac12 \br{e_p(\gen{v, w^\sigma}) + e_p(\gen{v, w^{\tau\sigma}})}.
    \]
    Applying the Cauchy--Schwarz inequality, we obtain
    \begin{align*}
        |\lambda_{v,w}|^2
        &\le \frac1{n!} \sum_{\sigma \in S_n}
        \frac14 \left|
            e_p(\gen{v, w^\sigma}) + e_p(\gen{v, w^{\tau\sigma}})
        \right|^2 \\
        &= \frac1{n!} \sum_{\sigma \in S_n}
        \frac14 \br{
            2 + 2 \Re\left[e_p(\gen{v, w^\sigma} - \gen{v, w^{\tau\sigma}})\right]
        }\\
        &= \frac12 + \frac12 \Re\left[\frac1{n!}\sum_{\sigma \in S_n}
            e_p(\langle v, (w - w^\tau)^\sigma\rangle)
        \right].
    \end{align*}
    Now $w - w^\tau = (u, -u, 0, 0, \dots)$, so
    \[
        \frac1{n!} \sum_{\sigma \in S_n} e_p(\gen{v, (w-w^\tau)^\sigma})
        = \frac1{n(n-1)} \sum_{\substack{1 \le i, j \le n \\ i \ne j}}
        e_p(u v_i - u v_j),
    \]
    and
    \[
        \sum_{\substack{1 \le i, j \le n \\ i \ne j}} e_p(u v_i - u v_j)
        = \sum_{1 \le i, j \le n} e_p(u v_i - u v_j) - n
        = n^2 |\lambda_v(u)|^2 - n,
    \]
    so we get
    \[
        |\lambda_{v,w}|^2
        \le \frac12 + \frac12 \frac{n|\lambda_v(u)|^2 - 1}{n-1}
        \le \frac12 + \frac12 |\lambda_v(u)|^2.
    \]
    This proves the lemma since $u \ne 0$.
\end{proof}

\end{document}
